\clearpage
\section{Local Schur comparison on the form
domain}\label{i-local-schur-comparison-on-the-form-domain}

Technical derivation for review. Equation and subsection numbers in this
appendix are local to the source derivation. Historical local labels are
retained, including numbering gaps or nonsequential labels. The
accompanying source notes retain the original expressions for
comparison.

\subsection{1. Statement and inputs}\label{i-statement-and-inputs}

Fix a partition of a finite \(N\). Put

% DISPLAY i.1
% SOURCE r_* = min_e r_e,    W_p = sum_e r_e^(-p),
% SOURCE H_s = H_I + T_C,    H_exc = sum_e(H_fast,e-E_0,e).
\begin{gather*}
\begin{aligned}
& r_{*} = \min _{e} r_{e}, \quad  W_{p} = \sum _{e} r_{e}^{-p}, \\
& H_{s} = H_{I} + T_{C}, \quad  H_{\mathrm{exc}} = \sum _{e}(H_{\mathrm{fast},e}-E_{0,e}).
\end{aligned}\notag
\end{gather*}

The internal tube is \(\alpha _{a}+\alpha _{b} \le  \kappa  r_{ab}\).
The fast tube has \(\vert Y\vert  < \sigma  r_{*}\). There is no bound
on ratios of unrelated \(r_{e}\) and no internal spectral-gap
assumption. \(H_{I}\) is the complete interacting colored internal form.
\(T_{C}\) on the full fiber retains the intrinsic normal connection. On
the slow bundle, \(H_{s}\) denotes the center-plus-internal model
identified by internal radial transport with the \(A=0\) vacuum line.
Its unperturbed internal forms are the colored ones; no lower-rank
physical-sector Hardy theorem is used.

For every \(\epsilon >0\), there are positive \(\kappa , \sigma , R\)
and a smooth even normalized cutoff vacuum \(U_{\sigma}\) supported
strictly inside the fast tube such that, for \(r_{*}\ge R\),

% DISPLAY i.2
% SOURCE (1-epsilon) H_s[f] - epsilon W_2[f]
% SOURCE    <= F_sigma[f]
% SOURCE    <= (1+epsilon) H_s[f] + epsilon W_2[f],              (1.1)
\begin{gather*}
\begin{aligned}
& (1-\epsilon ) H_{s}[f] - \epsilon  W_{2}[f] \\
& \le  F_{\sigma}[f] \\
& \le  (1+\epsilon ) H_{s}[f] + \epsilon  W_{2}[f],
\end{aligned}\tag{1.1}
\end{gather*}

where the actual zero-energy variational Hamiltonian Schur form is

% DISPLAY i.3
% SOURCE F_sigma[f] = inf_{v in Q(H), U_sigma^*v=0}
% SOURCE                       H[U_sigma f+v].                 (1.2)
\begin{gather*}
\begin{aligned}
& F_{\sigma}[f] = \inf _{v \in  Q(H), U_{\sigma}^{*}v=0} \\
& H[U_{\sigma} f+v].
\end{aligned}\tag{1.2}
\end{gather*}

Here Q(H) denotes the CLOSED QUADRATIC-FORM DOMAIN of the Hamiltonian,
not its operator domain D(H). Orthogonality is fiberwise over the slow
variables. For every \(f\) in \(Q(H_{s})\), the infimum is attained by a
unique \(v_{f}\) in Q(H) with \(U_{\sigma}^{*}v_{f}=0\). The form is
closed on the slow model\textquotesingle s form domain, interpreted with
Dirichlet support on the selected branch. The constants depend on the
fixed partition and \(\epsilon\), never on \(r_{\max}/r_{*}\).

Inputs proved in the companion notes are:

\begin{enumerate}
\def\labelenumi{\arabic{enumi}.}
\item
  The exact intrinsic kinetic identity and normalized matrices from
  Appendix A.
\item
  The rootwise all-vector connection estimates in Appendix \(E\) and
  differential-curvature identities in Appendix \(G\).
\item
  The actual H1/H2 inventory, finite source graph grading, and operator
  identity

  % DISPLAY i.4
  % SOURCE ||W-S^*H_exc^(-1)S|| <= C_N kappa W_2,              (1.3)
  \begin{gather*}
  \begin{aligned}
  & \Vert W-S^{*}H_{\mathrm{exc}}^{-1}S\Vert  \le  C_{N} \kappa  W_{2},
  \end{aligned}\tag{1.3}
  \end{gather*}

  from Appendix \(D\) . Here \(S=H1 U\) is the actual source, not a
  transported replacement.
\item
  The graph-local complete-slow-form bound for
  \(Z=-H_{\mathrm{exc}}^{-1}S\):

  % DISPLAY i.5
  % SOURCE H_s[Zf] <= C R^(-3)H_s[f]+C kappa W_2[f]+C W_5[f], (1.4)
  \begin{gather*}
  \begin{aligned}
  & H_{s}[Zf] \le  C R^{-3}H_{s}[f]+C \kappa  W_{2}[f]+C W_{5}[f],
  \end{aligned}\tag{1.4}
  \end{gather*}

  including spectator-vacuum derivatives, from Appendix \(C\) and the
  actual-source continuation.
\item
  The colored complementary reserve. A fixed positive part of the actual
  square \(J[v]=\Vert A p_{s} v\Vert ^{2}\) can be retained:

  % DISPLAY i.6
  % SOURCE H[v] >= (1-delta)(H_s[v]+H_exc[v])+c_0 J[v]         (1.5)
  \begin{gather*}
  \begin{aligned}
  & H[v] \ge  (1-\delta )(H_{s}[v]+H_{\mathrm{exc}}[v])+c_{0} J[v]
  \end{aligned}\tag{1.5}
  \end{gather*}

  on the cutoff complement, after choosing the parameters and exterior
  radius. For example \(c_{0}=1/2\) is allowed. This is NOT another copy
  of the original full Gauss square.
\end{enumerate}

This note supplies the mixed and diagonal-remainder estimates needed to
assemble these inputs. It never inverts \(H_{I}\) or expands a resolvent
in internal momenta.

\subsection{2. Exact error-column factorization, including
spin}\label{i-exact-error-column-factorization-including-spin}

Use the common positive-fast/positive-slow convention

% DISPLAY i.7
% SOURCE H_kin = ||p_s u||^2+||p_F u||^2+||a u+A p_s u+s u||^2,
% SOURCE p_s=(g_C^(-1/2)p_C,p_I),
% SOURCE a=W^(1/2)M^(-*) (C+N_Y)^*p_F,
% SOURCE s=-W^(1/2)M^(-*)S_f.
\begin{gather*}
\begin{aligned}
& H_{\mathrm{kin}} = \Vert p_{s} u\Vert ^{2}+\Vert p_{F} u\Vert ^{2}+\Vert a u+A p_{s} u+s u\Vert ^{2}, \\
& p_{s}=(g_{C}^{-1/2}p_{C},p_{I}), \\
& a=W^{1/2}M^{-*} (C+N_{Y})^{*}p_{F}, \\
& s=-W^{1/2}M^{-*}S_{f}.
\end{aligned}\notag
\end{gather*}

The scalar density term is added separately. Set

% DISPLAY i.8
% SOURCE d_0=F^(-1)C^*p_F,  P=L-J_2,  Z_m=F^(-1)P^*,
% SOURCE T=(I+Z_m)^(-1),  n=F^(-1)N_Y^*p_F,
% SOURCE X=n-Z_m d_0,  s_0=-F^(-1)S_f,
% SOURCE q=X+s_0,  X_1=n-F^(-1)L^*d_0,  X_2=F^(-1)J_2^*d_0.
\begin{gather*}
\begin{aligned}
& d_{0}=F^{-1}C^{*}p_{F}, \quad  P=L-J_{2}, \quad  Z_{m}=F^{-1}P^{*}, \\
& T=(I+Z_{m})^{-1}, \quad  n=F^{-1}N_{Y}^{*}p_{F}, \\
& X=n-Z_{m} d_{0}, \quad  s_{0}=-F^{-1}S_{f}, \\
& q=X+s_{0}, \quad  X_{1}=n-F^{-1}L^{*}d_{0}, \quad  X_{2}=F^{-1}J_{2}^{*}d_{0}.
\end{aligned}\notag
\end{gather*}

The symbol \(Z_{m}\) is an inverse-matrix error, distinct from the
dressing \(Z\). Exactly,

% DISPLAY i.9
% SOURCE a+s = W^(1/2)(d_0+Tq),
\begin{gather*}
\begin{aligned}
& a+s = W^{1/2}(d_{0}+Tq),
\end{aligned}\notag
\end{gather*}

and, as ordered first-order forms,

% DISPLAY i.10
% SOURCE ||(a+s)u||^2-||d_0u||^2
% SOURCE   =2 Re<d_0u,(X_1+s_0)u> + 2 Re<d_0u,X_2u>
% SOURCE    -2 Re<Z_m^*d_0u,Tq u> +||Tq u||^2
% SOURCE    +||K^*(d_0+Tq)u||^2.                              (2.1)
\begin{gather*}
\begin{aligned}
& \Vert (a+s)u\Vert ^{2}-\Vert d_{0} u\Vert ^{2} \\
& =2 \operatorname{Re} \langle d_{0} u,(X_{1}+s_{0})u\rangle  + 2 \operatorname{Re} \langle d_{0} u,X_{2} u\rangle \\
& -2 \operatorname{Re} \langle Z_{m}^{*}d_{0} u,Tq u\rangle  +\Vert Tq u\Vert ^{2} \\
& +\Vert K^{*}(d_{0}+Tq)u\Vert ^{2}.
\end{aligned}\tag{2.1}
\end{gather*}

The first line\textquotesingle s first term is exactly the
fast-orbital/spin part of H1. No derivative has been moved through an
inverse matrix. The remaining apparently dangerous \(d_{0}-X_{2}\)
product factors exactly:

% DISPLAY i.11
% SOURCE <d_0u,X_2u>
% SOURCE    =<C_Y F^(-1)d_0u,C_Y R^(-1)d_0u>.                 (2.2)
\begin{gather*}
\begin{aligned}
& \langle d_{0} u,X_{2} u\rangle \\
& =\langle C_{Y} F^{-1}d_{0} u,C_{Y} R^{-1}d_{0} u\rangle .
\end{aligned}\tag{2.2}
\end{gather*}

Every factor on the right of (2.1)-(2.2), after removing the displayed
H1 term, belongs to the finite error-column list

% DISPLAY i.12
% SOURCE X_1, X_2, Z_m^*d_0, K^*d_0,
% SOURCE C_Y F^(-1)d_0, C_Y R^(-1)d_0, s_0,
\begin{gather*}
\begin{aligned}
& X_{1}, X_{2}, Z_{m}^{*}d_{0}, K^{*}d_{0}, \\
& C_{Y} F^{-1}d_{0}, C_{Y} R^{-1}d_{0}, s_{0},
\end{aligned}\notag
\end{gather*}

followed by uniformly bounded left multiplication matrices. Their
all-vector estimates are

% DISPLAY i.13
% SOURCE ||E u||^2 <= C[(lambda H_exc)[u]+W_2[u]],
% SOURCE lambda=sigma^2+r_*^(-3).                              (2.3)
\begin{gather*}
\begin{aligned}
& \Vert E u\Vert ^{2} \le  C[(\lambda  H_{\mathrm{exc}})[u]+W_{2}[u]], \\
& \lambda =\sigma ^{2}+r_{*}^{-3}.
\end{aligned}\tag{2.3}
\end{gather*}

For the spin column the stronger bound
\(\Vert s_{0}u\Vert ^{2}\le CW_{2}[u]\) holds. On the uncut vacuum, and
on each fixed-degree Gaussian/Fock coefficient vector normalized in its
own roots,

% DISPLAY i.14
% SOURCE ||E U f||^2 <= C W_2[f].                              (2.4)
\begin{gather*}
\begin{aligned}
& \Vert E U f\Vert ^{2} \le  C W_{2}[f].
\end{aligned}\tag{2.4}
\end{gather*}

For (2.4), a same-edge \(Y_{e} p_{e}/r_{e}\) costs \(C/r_{e}\); a
triangle \(Y_{f} p_{e}/r_{g}\) costs \(C \sqrt{r_{e}/r_{f}}/r_{g}\),
whose square is bounded by \(C W_{2,t}\) by \(r_{e}\le r_{f}+r_{g}\).
Root-preserving pair matrices have uniformly bounded normalized
coefficients. All inverse matrices in the exact expressions remain
bounded LEFT factors. The same proof therefore works on a fixed finite
excited degree, with a degree-dependent constant, and does not introduce
\(r_{\max}\).

For a complementary \(v, W_{2}[v]\le C R^{-3}H_{\mathrm{exc}}[v]\).
Polarization of each product in (2.1)-(2.2) consequently gives, for any
\(\eta >0\),

% DISPLAY i.15
% SOURCE |R_fast[Uf,v]| <= eta H_exc[v]
% SOURCE                      +C_eta(sigma^2+R^(-3)) W_2[f].   (2.5)
\begin{gather*}
\begin{aligned}
& \vert R_{\mathrm{fast}}[Uf,v]\vert  \le  \eta  H_{\mathrm{exc}}[v] \\
& +C_{\eta}(\sigma ^{2}+R^{-3}) W_{2}[f].
\end{aligned}\tag{2.5}
\end{gather*}

This includes spin-fast and spin-square remainders. It is the central
reason a finite inverse-square constant is not left unpaid: no Taylor
remainder is ever paired directly with an unassigned \(d_{0}\) vacuum
norm.

For diagonal matching, retain the quadratic terms in (2.1), replace
\(T\) by I, \(q\) by \(X_{1}+s_{0}, Z_{m}\) by \(F^{-1}L^{*}\), and
retain (2.2) and \(\Vert K^{*}d_{0}\Vert ^{2}\). These are precisely the
corresponding H2 terms. Each discarded product has one additional
uniformly \(O(\sigma )\) factor and two low error columns. Hence

% DISPLAY i.16
% SOURCE |R_fast[Uf]-R_fast,H2[Uf]| <= C sigma W_2[f].          (2.6)
\begin{gather*}
\begin{aligned}
& \vert R_{\mathrm{fast}}[Uf]-R_{\mathrm{fast},H2}[Uf]\vert  \le  C \sigma  W_{2}[f].
\end{aligned}\tag{2.6}
\end{gather*}

One may obtain an additional inverse-radius improvement from Gaussian
moments, but it is unnecessary. The order of the two operations matters:
first perform the exact error-column factorization, then bound the
inverse difference. Doing them in reverse would expose the forbidden
remote vacuum mass.

\subsection{3. Other nondifferential fast
terms}\label{i-other-nondifferential-fast-terms}

The actual \(V_{3}\) and linear-Y Yukawa are exactly the remaining H1
terms; there is no omitted higher polynomial potential. They are
boson-odd, so their diagonal expectation on \(U\) or an even cutoff
\(U_{\sigma}\) is zero.

The quartic potential is a sum of squares of quadratic fast-coordinate
contractions. In addition to the distinct-root Q/P estimate already
proved, the repeated-root estimate is

% DISPLAY i.17
% SOURCE ||chi Y_e^2 u||^2 <= C sigma^2 h_e[u]+C r_e^(-2)||u||^2. (3.1)
\begin{gather*}
\begin{aligned}
& \Vert \chi  Y_{e}^{2} u\Vert ^{2} \le  C \sigma ^{2} h_{e}[u]+C r_{e}^{-2}\Vert u\Vert ^{2}.
\end{aligned}\tag{3.1}
\end{gather*}

Split the input into \(Q_{e}\) and \(P_{e}\), bound one \(Y_{e}\) by
\(\sigma  r_{*}\) on \(Q_{e}\), use
\(\Vert Y_{e} Q_{e} u\Vert ^{2}\le Cr_{e}^{-2}h_{e}[u]\), and use the
exact fourth Gaussian moment on \(P_{e}\). This gives

% DISPLAY i.18
% SOURCE 0<=V4[u]<=C sigma^2 H_exc[u]+C W_2[u],
% SOURCE V4[Uf]<=C W_2[f].                                    (3.2)
\begin{gather*}
\begin{aligned}
& 0\le V_{4}[u]\le C \sigma ^{2} H_{\mathrm{exc}}[u]+C W_{2}[u], \\
& V_{4}[Uf]\le C W_{2}[f].
\end{aligned}\tag{3.2}
\end{gather*}

Thus its low/complement cross obeys the same bound as (2.5). Its full
diagonal expectation belongs to \(W\) and is retained there.

The density multiplier satisfies \(\vert V_{d}\vert \le CW_{2}\). Its
mixed term is paid by the complementary L2 gap:

% DISPLAY i.19
% SOURCE |<Uf,V_d v>|<=eta H_exc[v]+C_eta R^(-3)W_2[f].         (3.3)
\begin{gather*}
\begin{aligned}
& \vert \langle Uf,V_{d} v\rangle \vert \le \eta  H_{\mathrm{exc}}[v]+C_{\eta} R^{-3}W_{2}[f].
\end{aligned}\tag{3.3}
\end{gather*}

The leading coefficient \(V_{d,2}\) is the one specified by the exact H2
inventory, including the slow density, fast Hessian, and orbital
ordering terms. Differentiating the normalized matrices in the exact
density formula and subtracting their zero-fast-coordinate jets gives

% DISPLAY i.20
% SOURCE |V_d-V_d,2|<=C sigma W_2                              (3.4)
\begin{gather*}
\begin{aligned}
& \vert V_{d}-V_{d,2}\vert \le C \sigma  W_{2}
\end{aligned}\tag{3.4}
\end{gather*}

on the tube. Each normalized inverse and its required derivatives are
bounded in its own root scale. More explicitly, \(V_{d}\) uses two
coefficient derivatives, and one additional fast derivative of that
finite rational expression is bounded by \(C/r_{*}^{3}\); integrating it
from \(Y=0\) gives \(C\vert Y\vert /r_{*}^{3}\le C \sigma  W_{2}\). The
leading zero-fast-coordinate coefficient includes the quadratic
determinant jet before fast differentiation. This is a
density-coefficient estimate, not a bound on \(d_{0}^{2}\). Equations
(3.3)-(3.4) settle both density uses.

The frozen \(E_{0}\) is a separate scalar multiplier. It has zero U-v
cross because it commutes with the vacuum projection. Its incident
payment is

% DISPLAY i.21
% SOURCE |E_0[u]|<=C R^(-3)H_I[u]+C kappa W_2[u].              (3.5)
\begin{gather*}
\begin{aligned}
& \vert E_{0}[u]\vert \le C R^{-3}H_{I}[u]+C \kappa  W_{2}[u].
\end{aligned}\tag{3.5}
\end{gather*}

It is never included in \(H_{\mathrm{exc}}\) or in a virtual
denominator.

\subsection{\texorpdfstring{4. Scalar fast-slow orbit cross: complete
charges instead of bare
\(t_{I}\)}{4. Scalar fast-slow orbit cross: complete charges instead of bare t\_\{I\}}}\label{i-scalar-fast-slow-orbit-cross-complete-charges-instead-of-bare-t_i}

Write \(\beta =A^{*}a\), with internal components
\(\beta _{j}=b_{j}(Y;s) \cdot  p_{F}\). The \(b_{j}\) are real, scalar
in internal Clifford factors, and commute with all internal
potential-charge coefficients. Let

% DISPLAY i.22
% SOURCE d_j=div_F b_j,    beta_hat,j=beta_j-(i/2)d_j,
% SOURCE C_alpha(beta_hat)=sum_j c_(alpha,j) beta_hat,j.
\begin{gather*}
\begin{aligned}
& d_{j}=\operatorname{div} _{F} b_{j}, \quad  \widehat{\beta}_{j}=\beta _{j}-(i/2)d_{j}, \\
& C_{\alpha}(\widehat{\beta})=\sum _{j} c_{\alpha ,j} \widehat{\beta}_{j}.
\end{aligned}\notag
\end{gather*}

The differentiated connection lemma and its explicit rootwise
decomposition imply

% DISPLAY i.23
% SOURCE sum_j ||beta_hat,j v||^2 <= C(lambda H_exc[v]+W_2[v]),
% SOURCE sum_j ||beta_hat,j U f||^2 <= C(kappa^2+sigma^2)W_2[f], (4.1)
\begin{gather*}
\begin{aligned}
& \sum _{j} \Vert \widehat{\beta}_{j} v\Vert ^{2} \le  C(\lambda  H_{\mathrm{exc}}[v]+W_{2}[v]), \\
& \sum _{j} \Vert \widehat{\beta}_{j} U f\Vert ^{2} \le  C(\kappa ^{2}+\sigma ^{2})W_{2}[f],
\end{aligned}\tag{4.1}
\end{gather*}

and

% DISPLAY i.24
% SOURCE sum_jk ||(partial_Aj beta_hat,k)U f||^2
% SOURCE                        <= C r_*^(-2)W_2[f].          (4.2)
\begin{gather*}
\begin{aligned}
& \sum _{jk} \Vert (\partial _{Aj} \widehat{\beta}_{k})U f\Vert ^{2} \\
& \le  C r_{*}^{-2}W_{2}[f].
\end{aligned}\tag{4.2}
\end{gather*}

The SMALL low-leg divergence estimate in (4.1) is essential. Directly,
\(b=A^{*}B\) with
\(A=O(\sigma ), B=O(\kappa +\sigma ), \partial _{Y} A=O(r_{*}^{-1})\)
and \(\partial _{Y} B=O(r_{*}^{-1})\). Hence

% DISPLAY i.25
% SOURCE |div_F b|<=C(kappa+sigma)/r_*.
\begin{gather*}
\begin{aligned}
& \vert \operatorname{div} _{F} b\vert \le C(\kappa +\sigma )/r_{*}.
\end{aligned}\notag
\end{gather*}

Together with the stronger rootwise estimate for \(A^{*}d_{0}\) and the
\(O(\sigma )\) factor multiplying \(a-d_{0}\), this proves the second
line of (4.1). It is not inferred from a coarse \(O(1/r_{*})\)
divergence bound.

\subsubsection{4.1 Exact linear charge
identity}\label{i-exact-linear-charge-identity}

Polarize the following identity, initially on the smooth core:

% DISPLAY i.26
% SOURCE 2 Re sum_j <p_j u,beta_j u>
% SOURCE   =2 Re (1/16)sum_alpha <Q_alpha u,C_alpha(beta_hat)u>
% SOURCE      -R_der[u]+(1/2)sum_j <u,(partial_Aj d_j)u>.      (4.3)
\begin{gather*}
\begin{aligned}
& 2 \operatorname{Re}  \sum _{j} \langle p_{j} u,\beta _{j} u\rangle \\
& =2 \operatorname{Re}  (1/16)\sum _{\alpha} \langle Q_{\alpha} u,C_{\alpha}(\widehat{\beta})u\rangle \\
& -R_{\mathrm{der}}[u]+(1/2)\sum _{j} \langle u,(\partial _{Aj} d_{j})u\rangle .
\end{aligned}\tag{4.3}
\end{gather*}

Here \(R_{\mathrm{der}}\) is the fixed bounded-Clifford contraction of
\(\partial _{j} \widehat{\beta}_{k}-\partial _{k} \widehat{\beta}_{j}\).
To verify it, expand the exact shifted averaged-charge identity and
subtract the beta-square. The beta-beta curvature in
\(C_{\alpha}(\widehat{\beta})^{2}\) cancels the beta-beta curvature in
the full shifted charge identity. Only the derivative curvature remains.
The raw-versus-symmetric cross contributes \(+(1/2)\partial _{A} d\),
with this sign. No anticommutator with the internal potential survives
the spinor average.

The first term in the polarized identity contains

% DISPLAY i.27
% SOURCE <Q(Uf),C(beta_hat)v>  and  <C(beta_hat)Uf,Qv>.
\begin{gather*}
\begin{aligned}
& \langle Q(Uf),C(\widehat{\beta})v\rangle  \quad  \text{ and } \quad  \langle C(\widehat{\beta})Uf,Qv\rangle .
\end{aligned}\notag
\end{gather*}

The first is bounded using (4.1), the complementary gap, and

% DISPLAY i.28
% SOURCE (1/16)sum ||Q_alpha(Uf)||^2 <= C(H_s[f]+W_2[f]).
\begin{gather*}
\begin{aligned}
& (1/16)\sum  \Vert Q_{\alpha}(Uf)\Vert ^{2} \le  C(H_{s}[f]+W_{2}[f]).
\end{aligned}\notag
\end{gather*}

The second uses the full \(H_{I}[v]\) reserve and the small low-leg
estimate in (4.1). No constant fraction of bare \(t_{I}\) is spent.
Since each derivative-curvature operator is symmetric, it can be moved
onto \(U\) in a mixed pairing. Equation (4.2) and the complementary gap
then give a \(C_{\eta} R^{-3}W_{2}\) cost. The scalar term has the same
payment.

Consequently the complete internal scalar-orbit mixed form obeys

% DISPLAY i.29
% SOURCE |R_orb,I[Uf,v]| <= eta(H_I[v]+H_exc[v])
% SOURCE    +C_eta lambda H_s[f]
% SOURCE    +C_eta(kappa^2+sigma^2+R^(-3))W_2[f].              (4.4)
\begin{gather*}
\begin{aligned}
& \vert R_{\mathrm{orb},I}[Uf,v]\vert  \le  \eta (H_{I}[v]+H_{\mathrm{exc}}[v]) \\
& +C_{\eta} \lambda  H_{s}[f] \\
& +C_{\eta}(\kappa ^{2}+\sigma ^{2}+R^{-3})W_{2}[f].
\end{aligned}\tag{4.4}
\end{gather*}

This is uniform over unrestricted internal momenta and colored internal
states. The center component obeys the same inequality by ordinary
center-kinetic Cauchy-Schwarz, with its actual normal connection
retained. Normal-frame generators are not estimated by a vacuum Berry
scalar on arbitrary \(v\).

\subsubsection{4.2 Independent incident-row
verification}\label{i-independent-incident-row-verification}

There is a second direct check of the high-momentum issue. With \(U=g\)
times its fermion vacuum and real Gaussian \(g\),

% DISPLAY i.30
% SOURCE (beta_j+beta_j^*)U
% SOURCE     =i(2 b_j dot Omega Y-div_F b_j)U.
\begin{gather*}
\begin{aligned}
& (\beta _{j}+\beta _{j}^{*})U \\
& =i(2 b_{j} \cdot  \Omega  Y-\operatorname{div} _{F} b_{j})U.
\end{aligned}\notag
\end{gather*}

Ground-state integration by parts therefore gives

% DISPLAY i.31
% SOURCE ||H_exc^(-1/2)(beta_j+beta_j^*)U||^2
% SOURCE                           <=C||b_jU||^2.            (4.5)
\begin{gather*}
\begin{aligned}
& \Vert H_{\mathrm{exc}}^{-1/2}(\beta _{j}+\beta _{j}^{*})U\Vert ^{2} \\
& \le C\Vert b_{j}U\Vert ^{2}.
\end{aligned}\tag{4.5}
\end{gather*}

For an internal row belonging to block a, factor the exact coefficient
as

% DISPLAY i.32
% SOURCE b_a=I_Y,a F^(-1)[T^* W M^(-*)(C+N_Y)^*].
\begin{gather*}
\begin{aligned}
& b_{a}=I_{Y,a} F^{-1}[T^{*} W M^{-*}(C+N_{Y})^{*}].
\end{aligned}\notag
\end{gather*}

The bracket is uniformly \(O(\kappa +\sigma )\), while
\(I_{Y,a} F^{-1}\) contains only \(Y_{e}/r_{e}\) on edges incident to a.
Thus

% DISPLAY i.33
% SOURCE ||b_aU||^2<=C(kappa+sigma)^2 w_a,
% SOURCE w_a=sum_(e incident a)r_e^(-3).                       (4.6)
\begin{gather*}
\begin{aligned}
& \Vert b_{a}U\Vert ^{2}\le C(\kappa +\sigma )^{2} w_{a}, \\
& w_{a}=\sum _{e \text{ incident } a}r_{e}^{-3}.
\end{aligned}\tag{4.6}
\end{gather*}

This proves directly that integrating the single slow derivative in the
orbit cross cannot create an unrelated internal-radius loss. It also
accounts for the adjoint fast divergence, rather than dropping it.
Either (4.3) or (4.5)-(4.6) supplies the needed first-derivative
estimate; the proof above uses (4.3).

\subsubsection{4.3 Diagonal orbit
remainder}\label{i-diagonal-orbit-remainder}

Let \(\beta _{1}\) be the first homogeneous slow-orbit coefficient: use
the linear slow column with \(F^{-1}\), the leading \(C_{Y}^{*}+FK\)
center term, and \(d_{0}\). Its raw diagonal compression, including
derivatives of its coefficient at \(A=0\), is retained in \(W\).
Together with the corresponding \(V_{d,2}\) orbital-density term, it
equals the \(V_{\mathrm{orb}}\) in the density-shifted inventory. In
particular the inventory term -2 \(\sum _{\mu} \rho _{0,\mu} j_{\mu}\)
belongs to \(V_{d,2}\) in the present unshifted-kinetic convention; it
is counted exactly once, not inserted a second time in the raw
\(\beta _{1}\) compression. Only the SUM of these contributions is
identified with the H2 scalar. The raw \(\beta _{1}\) contribution is
NOT discarded on the grounds that \(\beta _{1}\) vanishes at \(A=0\).

The exact identities for \(A^{*}d_{0}\) in the rootwise helper express
\(\beta -\beta _{1}\) as a sum with one extra \(O(\sigma )\) factor
multiplying a controlled error column. They give

% DISPLAY i.34
% SOURCE ||(beta-beta_1)U||^2+||(div(b-b_1))U||^2<=C sigma^2 W_2,
% SOURCE ||partial_A(beta-beta_1)U||<=C sigma r_*^(-1) W_2^(1/2).
\begin{gather*}
\begin{aligned}
& \Vert (\beta -\beta _{1})U\Vert ^{2}+\Vert (\operatorname{div} (b-b_{1}))U\Vert ^{2}\le C \sigma ^{2} W_{2}, \\
& \Vert \partial _{A}(\beta -\beta _{1})U\Vert \le C \sigma  r_{*}^{-1} W_{2}^{1/2}.
\end{aligned}\notag
\end{gather*}

Differentiate those identities, not an abstract norm bound: the extra
fast-degree factor remains when \(F\) or \(F^{-1}C^{*}\) is
differentiated. Applying (4.3) to this difference gives, for any
\(\eta >0\),

% DISPLAY i.35
% SOURCE |R_orb,diag[f]|<=eta H_s[f]
% SOURCE                          +C_eta(sigma+sigma^2)W_2[f]. (4.7)
\begin{gather*}
\begin{aligned}
& \vert R_{\mathrm{orb},\mathrm{diag}}[f]\vert \le \eta  H_{s}[f] \\
& +C_{\eta}(\sigma +\sigma ^{2})W_{2}[f].
\end{aligned}\tag{4.7}
\end{gather*}

This retains the exact leading coefficient while paying only its true
higher-order remainder.

\subsection{5. Spin-slow, slow-orbit square, and induced center
metric}\label{i-spin-slow-slow-orbit-square-and-induced-center-metric}

\subsubsection{5.1 Spin-slow cross}\label{i-spin-slow-cross}

Put \(c=A^{*}s\). Its internal row a satisfies the exact incident
coefficient bound

% DISPLAY i.36
% SOURCE |c_a|^2<=C W_2 sum_(e incident a)|Y_e|^2/r_e^2,
% SOURCE ||c_a U||^2<=C w_a W_2.                              (5.1)
\begin{gather*}
\begin{aligned}
& \vert c_{a}\vert ^{2}\le C W_{2} \sum _{e \text{ incident } a}\vert Y_{e}\vert ^{2}/r_{e}^{2}, \\
& \Vert c_{a} U\Vert ^{2}\le C w_{a} W_{2}.
\end{aligned}\tag{5.1}
\end{gather*}

For the mixed form, integrate its single slow derivative onto the smooth
coefficient cU \(f\). This is legitimate for the complete matrix \(c\),
with no commutation claim about internal Clifford matrices. The
coefficient of \(\partial _{Aa} f\) has norm squared bounded by
\(C w_{a} W_{2}\). The complementary gap therefore gives the weighted
derivative cost

% DISPLAY i.37
% SOURCE C(W_2/r_*) sum_a w_a t_a[f]
% SOURCE    <=C R^(-3)[R^(-3)H_I[f]+kappa W_2[f]].             (5.2)
\begin{gather*}
\begin{aligned}
& C(W_{2}/r_{*}) \sum _{a} w_{a} t_{a}[f] \\
& \le C R^{-3}[R^{-3}H_{I}[f]+\kappa  W_{2}[f]].
\end{aligned}\tag{5.2}
\end{gather*}

The exact inverse derivative formula, incident left factor
\(I_{Y,a} F^{-1}\), and finite Gaussian derivative moments give

% DISPLAY i.38
% SOURCE sum_a ||partial_Aa(c_aU)||^2
% SOURCE    <=C r_*^(-2)(sum_a w_a)W_2+C(sum_a w_a)W_2^2
% SOURCE    <=C R^(-5)W_2.                                   (5.3)
\begin{gather*}
\begin{aligned}
& \sum _{a} \Vert \partial _{Aa}(c_{a} U)\Vert ^{2} \\
& \le C r_{*}^{-2}(\sum _{a} w_{a})W_{2}+C(\sum _{a} w_{a})W_{2}^{2} \\
& \le C R^{-5}W_{2}.
\end{aligned}\tag{5.3}
\end{gather*}

Spectator-vacuum derivatives are included in the \(W_{2}^{2}\) term.
Center derivatives have the same or better bound and are paid by
\(T_{C}\). After the complementary gap, (5.2)-(5.3) are smaller than the
mixed budget in (4.4).

For the DIAGONAL, (5.1) alone would leave a fixed \(W_{2}\) constant and
is insufficient. The leading coefficient \(c_{1}\) is boson-odd: its
internal part is \(-I_{Y} F^{-2}S_{f}\), with the corresponding linear
center row. \(U\), its slow derivatives, and the even cutoff are
boson-even. Its diagonal cross is exactly zero. Every term of
\(c-c_{1}\) has an additional inverse/metric difference of size
\(O(\sigma )\), so

% DISPLAY i.39
% SOURCE ||(c_a-c_1,a)U||^2<=C sigma^2 w_a W_2.
\begin{gather*}
\begin{aligned}
& \Vert (c_{a}-c_{1,a})U\Vert ^{2}\le C \sigma ^{2} w_{a} W_{2}.
\end{aligned}\notag
\end{gather*}

Weighted Young and \(\sum _{a} w_{a} \alpha _{a}\le C \kappa  W_{2}\)
yield only

% DISPLAY i.40
% SOURCE eta R^(-3)H_s[f]+C_eta(kappa+sigma^2)W_2[f].          (5.4)
\begin{gather*}
\begin{aligned}
& \eta  R^{-3}H_{s}[f]+C_{\eta}(\kappa +\sigma ^{2})W_{2}[f].
\end{aligned}\tag{5.4}
\end{gather*}

This is the specific parity use that prevents an unpaid diagonal
constant.

\subsubsection{5.2 The positive square actually retained by
completion}\label{i-the-positive-square-actually-retained-by-completion}

Let \(J[u]=\Vert A p_{s} u\Vert ^{2}\). Its low-leg estimate is

% DISPLAY i.41
% SOURCE J[Uf]<=C R^(-3)H_s[f]+C kappa W_2[f]+C W_5[f].        (5.5)
\begin{gather*}
\begin{aligned}
& J[Uf]\le C R^{-3}H_{s}[f]+C \kappa  W_{2}[f]+C W_{5}[f].
\end{aligned}\tag{5.5}
\end{gather*}

For derivatives on \(f\), the exact incident factorization gives
\(\sum _{a} w_{a} t_{a}\) and the analogous center coefficient. Convert
only these incident weights using

% DISPLAY i.42
% SOURCE sum_a w_a t_a<=C R^(-3)H_I+C kappa W_2.
\begin{gather*}
\begin{aligned}
& \sum _{a} w_{a} t_{a}\le C R^{-3}H_{I}+C \kappa  W_{2}.
\end{aligned}\notag
\end{gather*}

Derivatives on \(U\) are finite degree-two Gaussian/Fock vectors. Their
rootwise moments give \(W_{5}\), including spectator products
\(W_{3}W_{2}\le C_{N}W_{5}\). The Berry connection adds its already
proved smaller incident contribution. Thus

% DISPLAY i.43
% SOURCE |<A p_sUf,A p_sv>|
% SOURCE    <=eta J[v]+C_eta[R^(-3)H_s[f]+kappa W_2[f]+W_5[f]]. (5.6)
\begin{gather*}
\begin{aligned}
& \vert \langle A p_{s} U f,A p_{s} v\rangle \vert \\
& \le \eta  J[v]+C_{\eta}[R^{-3}H_{s}[f]+\kappa  W_{2}[f]+W_{5}[f]].
\end{aligned}\tag{5.6}
\end{gather*}

Equation (5.6) uses only the true postcompletion reserve \(J\). No
original full Gauss square is used.

\subsubsection{5.3 Center metric}\label{i-center-metric}

The direct center coefficient differs from identity by

% DISPLAY i.44
% SOURCE g_C^(-1)-I=-K^*(I+KK^*)^(-1)K.
\begin{gather*}
\begin{aligned}
& g_{C}^{-1}-I=-K^{*}(I+KK^{*})^{-1}K.
\end{aligned}\notag
\end{gather*}

It is \(O(\sigma ^{2})\) on arbitrary tube vectors. On a Gaussian low
leg its squared coefficient norm is \(O(R^{-6})\); derivatives on \(U\)
have the corresponding \(W_{5}\)-or-smaller moments. Ordinary
center-form Young controls the mixed term by
\(\eta  T_{C}[v]+C_{\eta} R^{-6}T_{C}[f]+C_{\eta}W_{5}[f]\). Its
diagonal has bound \(C R^{-3}T_{C}[f]+CW_{5}[f]\). No internal
derivative is involved.

\subsubsection{5.4 Direct slow kinetic
cross}\label{i-direct-slow-kinetic-cross}

The preceding orbit estimates do not replace the cross from the DIRECT
center and internal kinetic forms. For \(v\) perpendicular to \(U\),
differentiate \(U^{*}v=0\) and integrate the remaining slow derivative
once. In a covariant vacuum-line frame the complementary source is

% DISPLAY i.45
% SOURCE -2 sum_j (1-P)(nabla_j U) nabla_j f
% SOURCE       -(1-P)sum_j nabla_j^2 U f.                    (5.7)
\begin{gather*}
\begin{aligned}
& -2 \sum _{j} (1-P)(\nabla _{j} U) \nabla _{j} f \\
& -(1-P)\sum _{j} \nabla _{j}^{2} U f.
\end{aligned}\tag{5.7}
\end{gather*}

This is a form identity and contains no second derivative of \(f\).
Internal potential and Yukawa terms commute with the vacuum injection
and have zero mixed compression. The same is true of \(E_{0}\).
Normal-bundle derivatives are kept covariantly in (5.7).

An internal derivative in block a differentiates only incident pair
vacua. After removing its scalar Berry part, each term has a fixed
excited sector on that pair, norm \(C/r_{e}\), and gap at least
\(c r_{e}\). Its squared \(H_{\mathrm{exc}}\)-dual norm is therefore at
most \(C w_{a}\). The center derivative satisfies the analogous sum of
pair weights. Second derivatives have active-pair and two-pair finite
excitation terms, of norm at most \(C/r_{e}^{2}\) or
\(C/(r_{e} r_{f})\). Their dual squared norm is bounded by \(C W_{5}\);
spectator products are included using \(W_{3}W_{2}\le C_{N}W_{5}\). Thus
ordinary source/gap Young gives

% DISPLAY i.46
% SOURCE 2|H_s[Uf,v]| <= eta H_exc[v]
% SOURCE        +C_eta[R^(-3)H_s[f]+kappa W_2[f]+W_5[f]].    (5.8)
\begin{gather*}
\begin{aligned}
& 2\vert H_{s}[Uf,v]\vert  \le  \eta  H_{\mathrm{exc}}[v] \\
& +C_{\eta}[R^{-3}H_{s}[f]+\kappa  W_{2}[f]+W_{5}[f]].
\end{aligned}\tag{5.8}
\end{gather*}

The conversion on the right is incident:
\(\sum _{a} w_{a} t_{a}\le C R^{-3}H_{I}+C \kappa  W_{2}\). It is
performed before replacing \(w_{a}\) by its maximum. The cutoff version
gains only the source-specific exponentially small errors in Section
\(7\). This explicitly includes the direct slow derivative source in the
mixed ledger.

\subsection{6. Summary of the mixed and diagonal
estimates}\label{i-summary-of-the-mixed-and-diagonal-estimates}

Let \(B[v]=H_{s}[v]+H_{\mathrm{exc}}[v]+J[v]\). After a fixed finite
allocation of \(\eta\) to the preceding terms, for every \(\eta >0\) the
exact mixed form has the decomposition

% DISPLAY i.47
% SOURCE H[Uf,v]=<Sf,v>+R_mix[f,v],
\begin{gather*}
\begin{aligned}
& H[Uf,v]=\langle Sf,v\rangle +R_{\mathrm{mix}}[f,v],
\end{aligned}\notag
\end{gather*}

with

% DISPLAY i.48
% SOURCE 2|R_mix[f,v]|
% SOURCE    <=eta B[v]+a_(eta)(kappa,sigma,R)H_s[f]
% SOURCE                      +b_(eta)(kappa,sigma,R)W_2[f],  (6.1)
\begin{gather*}
\begin{aligned}
& 2\vert R_{\mathrm{mix}}[f,v]\vert \\
& \le \eta  B[v]+a_{\eta}(\kappa ,\sigma ,R)H_{s}[f] \\
& +b_{\eta}(\kappa ,\sigma ,R)W_{2}[f],
\end{aligned}\tag{6.1}
\end{gather*}

where one can choose finite constants such that

% DISPLAY i.49
% SOURCE a_(eta)<=C_eta(sigma^2+R^(-3)),
% SOURCE b_(eta)<=C_eta(kappa+kappa^2+sigma^2+R^(-3)).          (6.2)
\begin{gather*}
\begin{aligned}
& a_{\eta}\le C_{\eta}(\sigma ^{2}+R^{-3}), \\
& b_{\eta}\le C_{\eta}(\kappa +\kappa ^{2}+\sigma ^{2}+R^{-3}).
\end{aligned}\tag{6.2}
\end{gather*}

The exact diagonal is

% DISPLAY i.50
% SOURCE H[Uf]=H_s^B[f]+E_0[f]+W[f]+R_diag[f],                (6.3)
\begin{gather*}
\begin{aligned}
& H[Uf]=H_{s}^{B}[f]+E_{0}[f]+W[f]+R_{\mathrm{diag}}[f],
\end{aligned}\tag{6.3}
\end{gather*}

and for an independently chosen \(\theta >0\),

% DISPLAY i.51
% SOURCE |R_diag[f]|<=theta H_s[f]
% SOURCE      +C_theta[kappa+sigma+sigma^2+R^(-3)]W_2[f]
% SOURCE      +C R^(-3)H_s[f].                               (6.4)
\begin{gather*}
\begin{aligned}
& \vert R_{\mathrm{diag}}[f]\vert \le \theta  H_{s}[f] \\
& +C_{\theta}[\kappa +\sigma +\sigma ^{2}+R^{-3}]W_{2}[f] \\
& +C R^{-3}H_{s}[f].
\end{aligned}\tag{6.4}
\end{gather*}

The \(H_{s}^{B}\) comparison is the full averaged-charge Berry
comparison from source Section \(4.2\) of Appendix \(F\):

% DISPLAY i.52
% SOURCE |H_s^B[f]-H_s[f]|
% SOURCE    <=theta H_s[f]+C(theta^(-1)kappa^4+kappa)W_2[f].  (6.5)
\begin{gather*}
\begin{aligned}
& \vert H_{s}^{B}[f]-H_{s}[f]\vert \\
& \le \theta  H_{s}[f]+C(\theta ^{-1}\kappa ^{4}+\kappa )W_{2}[f].
\end{aligned}\tag{6.5}
\end{gather*}

Equations (3.5), (6.3)-(6.5) keep every internal potential and Yukawa
term. No frozen slow-momentum resolvent has appeared.

\subsection{7. Cutoff vacuum, derivative tails, and source
parity}\label{i-cutoff-vacuum-derivative-tails-and-source-parity}

Fix \(\sigma >0\) before increasing \(R\). Let
\(h=(\sum _{e} r_{e}^{-2})^{-1/2}\), and choose a smooth even radial
\(\chi =\chi _{0}(\vert Y\vert ^{2}/(\sigma _{0}^{2}h^{2}))\), with
\(0<\sigma _{0}<c_{N} \sigma\) fixed, equal to one on an inner tube and
supported strictly inside the validity tube. Put

% DISPLAY i.53
% SOURCE U_sigma=chi U/n,    n=||chi U||.
\begin{gather*}
\begin{aligned}
& U_{\sigma}=\chi  U/n, \quad  n=\Vert \chi  U\Vert .
\end{aligned}\notag
\end{gather*}

This is intrinsic and residual-equivariant. Its support is
projection-stable. The normalized Gaussian coordinates
\(\xi _{e}=r_{e}^{1/2}Y_{e}\) have uniformly positive bounded covariance
matrices. On the cutoff transition region,

% DISPLAY i.54
% SOURCE sum_e |xi_e|^2 >= c sigma_0^2 r_*^3.
\begin{gather*}
\begin{aligned}
& \sum _{e} \vert \xi _{e}\vert ^{2} \ge  c \sigma _{0}^{2} r_{*}^{3}.
\end{aligned}\notag
\end{gather*}

Consequently any fixed-degree root-normalized Gaussian polynomial has a
tail bounded by

% DISPLAY i.55
% SOURCE tau_R=C_(N,sigma_0) exp(-c_N sigma_0^2 R^3),           (7.1)
\begin{gather*}
\begin{aligned}
& \tau _{R}=C_{N,\sigma _{0}} \exp (-c_{N} \sigma _{0}^{2} R^{3}),
\end{aligned}\tag{7.1}
\end{gather*}

after absorbing fixed powers of \(R\) into the exponential. These
estimates do not contain \(r_{\max}\). They apply after the rootwise
error-column factorization, and to finite graph source/dressing
coefficients with their own graph weights.

One must not claim that \(\Vert p_{F}(U_{\sigma}-U)\Vert\) has a
root-independent bound: a far vacuum derivative can indeed carry an
arbitrarily large zero-point norm. Instead use the excitation
ground-state transform. Since

% DISPLAY i.56
% SOURCE H_exc U=0,
% SOURCE H_exc[chi U]=integral |grad_F chi|_G^2 |U|^2,
\begin{gather*}
\begin{aligned}
& H_{\mathrm{exc}} U=0, \\
& H_{\mathrm{exc}}[\chi  U]=\int  \vert \nabla _{F} \chi \vert _{G}^{2} \vert U\vert ^{2},
\end{aligned}\notag
\end{gather*}

no spectator vacuum energy occurs. In the operator cutoff commutator,
\(\nabla _{e} \chi\) is a scalar multiple of \(Y_{e}/h^{2}\); its
product with \(\Omega _{e}Y_{e}\) has a root-normalized quadratic
coefficient of order \(h^{-2}\), because \(\Omega _{e}/r_{e}\) is
bounded. The same observation applies to the active finite graph states
in \(Z\) and to H1: replacing a \(p_{e}\) by \(\nabla _{e} \chi\) adds a
factor bounded by \(C/(r_{e} h^{2})\), with a Gaussian tail. Thus the
only cutoff errors required here satisfy, uniformly in all root ratios,

% DISPLAY i.57
% SOURCE |delta diagonal|<=tau_R(H_s[f]+W_2[f]),
% SOURCE |delta mixed|<=eta B[v]+C_eta tau_R(H_s[f]+W_2[f]),
% SOURCE H_exc[chi Zf/n]=H_exc[Zf]+O(tau_R W_2[f]),             (7.2)
\begin{gather*}
\begin{aligned}
& \vert \delta  \text{ diagonal }\vert \le \tau _{R}(H_{s}[f]+W_{2}[f]), \\
& \vert \delta  \text{ mixed }\vert \le \eta  B[v]+C_{\eta} \tau _{R}(H_{s}[f]+W_{2}[f]), \\
& H_{\mathrm{exc}}[\chi  Zf/n]=H_{\mathrm{exc}}[Zf]+O(\tau _{R} W_{2}[f]),
\end{aligned}\tag{7.2}
\end{gather*}

with the analogous graph-local \(H_{s}\) and \(J\) estimates. Slow
cutoff derivatives use \(\vert \partial  h\vert /h\le C/r_{*}\);
internal derivatives of \(h\) vanish. Derivatives of \(n\) are
controlled by the same normalized tail calculation. The normal
connection preserves \(\vert Y\vert ^{2}\), so it creates no extra
radial-cutoff generator.

\(S\) and \(Z\) are boson-odd. The cutoff is even, hence \(\chi\) Z/n is
EXACTLY orthogonal to \(U_{\sigma}\). The leading cutoff H1 source is
also odd and is orthogonal to the uncut \(U\), so its shifted-fast
inverse never encounters the vacuum kernel. Replacing it by \(S\) in the
mixed identity costs only (7.2). In the lower bound the source
completion with the uncut \(H_{\mathrm{exc}}\) is valid for every \(v\),
whether or not \(v\) is exactly orthogonal to \(U\): \(S\) is orthogonal
to its kernel.

Finally \(v\) perp \(U_{\sigma}\) has
\(\Vert Pv\Vert \le \tau _{R}\Vert v\Vert\) and therefore

% DISPLAY i.58
% SOURCE H_exc[v]>=c r_*(1-tau_R^2)||v||^2.                   (7.3)
\begin{gather*}
\begin{aligned}
& H_{\mathrm{exc}}[v]\ge c r_{*}(1-\tau _{R}^{2})\Vert v\Vert ^{2}.
\end{aligned}\tag{7.3}
\end{gather*}

All estimates (6.1)-(6.4) consequently hold for \(U_{\sigma}\) and its
genuine Dirichlet complement after adding the vanishing \(\tau _{R}\)
terms.

\subsection{8. Lower Schur comparison}\label{i-lower-schur-comparison}

Choose the reserve loss \(\delta\) and the mixed allocation \(\eta\)
small. A fixed positive fraction of \(J\) is available, so \(\eta\) can
be chosen below \(c_{0}\). Equations (1.5) and (6.1) leave

% DISPLAY i.59
% SOURCE H[U_sigma f+v]
% SOURCE   >=H[U_sigma f]+c_s H_s[v]+c_J J[v]
% SOURCE        +(1-delta-eta)H_exc[v]+2Re<Sf,v>
% SOURCE        -a_eta H_s[f]-b_eta W_2[f],                  (8.1)
\begin{gather*}
\begin{aligned}
& H[U_{\sigma} f+v] \\
& \ge H[U_{\sigma} f]+c_{s} H_{s}[v]+c_{J} J[v] \\
& +(1-\delta -\eta )H_{\mathrm{exc}}[v]+2\operatorname{Re} \langle Sf,v\rangle \\
& -a_{\eta} H_{s}[f]-b_{\eta} W_{2}[f],
\end{aligned}\tag{8.1}
\end{gather*}

where \(c_{s}\) and \(c_{J}\) are positive after the finite allocation.
Completing only the shifted fast form gives

% DISPLAY i.60
% SOURCE (1-delta-eta)H_exc[v]+2Re<Sf,v>
% SOURCE   >=-(1-delta-eta)^(-1)<f,S^*H_exc^(-1)Sf>.          (8.2)
\begin{gather*}
\begin{aligned}
& (1-\delta -\eta )H_{\mathrm{exc}}[v]+2\operatorname{Re} \langle Sf,v\rangle \\
& \ge -(1-\delta -\eta )^{-1}\langle f,S^{*}H_{\mathrm{exc}}^{-1}Sf\rangle .
\end{aligned}\tag{8.2}
\end{gather*}

The virtual multiplier is \(\le CW_{2}\). Therefore replacing its
coefficient by one costs only \(C(\delta +\eta )W_{2}\). Use (1.3),
(3.5), and (6.3)-(6.5). Their total error is made smaller than
\(\epsilon (H_{s}+W_{2})\) in the parameter order below. Taking the
infimum proves the lower half of (1.1). One may retain a further small
positive fraction of \(H_{s}[v]+H_{\mathrm{exc}}[v]+J[v]\) for the
form-domain argument, at an additional \(O(\eta  W_{2})\) cost.

\subsection{9. Upper trial and the complete internal-source graph
bound}\label{i-upper-trial-and-the-complete-internal-source-graph-bound}

Use the genuine complementary trial

% DISPLAY i.61
% SOURCE v_f=chi Zf/n,   Z=-H_exc^(-1)S.
\begin{gather*}
\begin{aligned}
& v_{f}=\chi  Zf/n, \quad  Z=-H_{\mathrm{exc}}^{-1}S.
\end{aligned}\notag
\end{gather*}

The actual graph components obey

% DISPLAY i.62
% SOURCE ||Z_e||^2<=C r_e^(-3),
% SOURCE ||Z_t||^2<=C w_t/s_t,
% SOURCE w_t=sum_(e in t)r_e^(-2),  s_t=sum_(e in t)r_e,
\begin{gather*}
\begin{aligned}
& \Vert Z_{e}\Vert ^{2}\le C r_{e}^{-3}, \\
& \Vert Z_{t}\Vert ^{2}\le C w_{t}/s_{t}, \\
& w_{t}=\sum _{e \in  t}r_{e}^{-2}, \quad  s_{t}=\sum _{e \in  t}r_{e},
\end{aligned}\notag
\end{gather*}

and their full derivatives include the spectator term \(m_{g}W_{2}\).
The complete interacting internal estimate is (1.4). At a graph vertex,
\(\alpha _{a} m_{g}\le C \kappa  W_{2,g}\); outside the graph the
potential charge commutes with the even graph map, and only
differentiated spectator vacua remain. This is why no \(\alpha _{a}\)
from an unrelated far block can multiply a nearest-root dressing norm.
The rootwise graph proof covers arbitrary momenta of \(f\).

There is also a direct incident estimate for the retained orbit square:

% DISPLAY i.63
% SOURCE J[Zf]<=C R^(-6)H_s[f]
% SOURCE                   +C kappa R^(-3)W_2[f]+C sigma^2 W_5[f]. (9.1)
\begin{gather*}
\begin{aligned}
& J[Zf]\le C R^{-6}H_{s}[f] \\
& +C \kappa  R^{-3}W_{2}[f]+C \sigma ^{2} W_{5}[f].
\end{aligned}\tag{9.1}
\end{gather*}

Indeed each fixed-degree graph state has the same \(Y_{e}\)
second-moment scale \(C/r_{e}\) on active and spectator roots. The
coefficient of \(t_{a}[f]\) is at most \(C w_{a} \sum _{g} m_{g}\).
Convert \(w_{a} t_{a}\) BEFORE using \(\sum _{g} m_{g}\le CR^{-3}\).
Derivatives on \(Z\) use the full active-plus-spectator derivative bound
and the pointwise coefficient \(\Vert A\Vert \le C \sigma\), giving
\(C \sigma ^{2} W_{5}\). This proves (9.1) without a global internal
kinetic loss.

The same absolute estimates used in the reserve proof, with the upper
Young inequality for the complete shifted charges, give

% DISPLAY i.64
% SOURCE H[w]<= (1+delta)(H_s[w]+H_exc[w])
% SOURCE                     +C_delta W_2[w]+C J[w]+V4[w].   (9.2)
\begin{gather*}
\begin{aligned}
& H[w]\le  (1+\delta )(H_{s}[w]+H_{\mathrm{exc}}[w]) \\
& +C_{\delta} W_{2}[w]+C J[w]+V_{4}[w].
\end{aligned}\tag{9.2}
\end{gather*}

For the fast part this is immediate from the exact identity (2.1), whose
leading signed terms have two-sided bounds. The internal shifted-charge
comparison is also two-sided; density, E0 and spin estimates are
absolute. If desired \(V_{4}\) can be absorbed using (3.2). Thus (9.2)
uses no unproved upper control of the full original Gauss square.

Apply (9.2) to \(v_{f}\). Since
\(H_{\mathrm{exc}}[Zf]=\langle f,S^{*}H_{\mathrm{exc}}^{-1}Sf\rangle\),
(1.4), (9.1), (3.2), and the cutoff estimates give

% DISPLAY i.65
% SOURCE H[v_f]<= (1+delta+C sigma^2)H_exc[Zf]
% SOURCE    +C R^(-3)H_s[f]+C kappa W_2[f]
% SOURCE    +C_delta R^(-3)W_2[f]+O(tau_R(H_s+W_2)).          (9.3)
\begin{gather*}
\begin{aligned}
& H[v_{f}]\le  (1+\delta +C \sigma ^{2})H_{\mathrm{exc}}[Zf] \\
& +C R^{-3}H_{s}[f]+C \kappa  W_{2}[f] \\
& +C_{\delta} R^{-3}W_{2}[f]+O(\tau _{R}(H_{s}+W_{2})).
\end{aligned}\tag{9.3}
\end{gather*}

The mixed remainder (6.1) applied to this trial is small by the same
graph estimates. Meanwhile

% DISPLAY i.66
% SOURCE 2Re<Sf,Zf>=-2<f,S^*H_exc^(-1)Sf>.
\begin{gather*}
\begin{aligned}
& 2\operatorname{Re} \langle Sf,Zf\rangle =-2\langle f,S^{*}H_{\mathrm{exc}}^{-1}Sf\rangle .
\end{aligned}\notag
\end{gather*}

Together with the diagonal expansion, this leaves
\(W-S^{*}H_{\mathrm{exc}}^{-1}S\) and only the displayed arbitrarily
small errors. Equation (1.3) therefore proves the upper half of (1.1).
The graph trial is used only for this upper bound; arbitrary
complementary vectors remain unrestricted in the lower bound.

\subsection{10. Simultaneous parameter
order}\label{i-simultaneous-parameter-order}

There is no circular order and no unattached constant \(C W_{2}\):

\begin{enumerate}
\def\labelenumi{\arabic{enumi}.}
\tightlist
\item
  Fix the partition and desired \(\epsilon\). Choose the finite Young
  allocations \(\theta , \eta\) and reserve loss \(\delta\) so the
  coefficients they alone contribute, including \(C(\delta +\eta )\)
  from (8.2), use less than epsilon/10 each.
\item
  Choose \(\kappa\) small enough for the fixed pair functional-calculus
  neighborhood, \(C_{\eta} \kappa , C \kappa\) from the source defect
  and E0 payment, and \(C \theta ^{-1}\kappa ^{4}\) to fit their
  assigned \(W_{2}\) budgets.
\item
  Choose a FIXED positive \(\sigma\) small enough for the reserve,
  \(C_{\theta} \sigma , C_{\eta} \sigma ^{2}\) and the complete-charge
  mixed costs to fit their budgets. Choose a fixed smaller
  \(\sigma _{0}\) for the radial cutoff.
\item
  Increase \(R\) until every constant now fixed times \(R^{-3}\), all
  complementary \(W_{2}\) absorptions, and all cutoff-tail quantities
  fit the remaining budgets.
\end{enumerate}

Constants such as \(C_{\delta}\) or \(1/\sigma\) are finite before \(R\)
is chosen. No estimate uses a later inverse power of a parameter to undo
an earlier smallness requirement without allowing that earlier parameter
to be chosen in the stated order. All estimates are uniform on nested
smaller supports at the selected fixed widths.

\subsection{11. Closed domains and variational
meaning}\label{i-closed-domains-and-variational-meaning}

First work on compact slow supports strictly inside the chosen branch
and tube. The exact first-order coefficients and density are smooth
there, with uniformly positive principal kinetic form and bounded local
matrix potentials. The associated Dirichlet form is closed after a
scalar shift, with its ordinary local H1 core. \(U_{\sigma}\) has smooth
compact fast support; \(P_{\sigma}\) and \(Q_{\sigma}\) map this core
into the same form domain. No assertion is made that the uncut \(P\)
preserves Dirichlet support.

The estimates above imply, after adding a bounded scalar multiple of
\(\Vert u\Vert ^{2}\),

% DISPLAY i.67
% SOURCE H[u]+C_R||u||^2 >= c H_s[f]+c B[v],
% SOURCE u=U_sigma f+v,
\begin{gather*}
\begin{aligned}
& H[u]+C_{R}\Vert u\Vert ^{2} \ge  c H_{s}[f]+c B[v], \\
& u=U_{\sigma} f+v,
\end{aligned}\notag
\end{gather*}

with \(c>0\); \(W_{2}\le C R^{-2}\) is bounded. Conversely the diagonal
upper estimate bounds \(H[U_{\sigma} f]\) by
\(C(H_{s}[f]+\Vert f\Vert ^{2})\). Hence \(P_{\sigma}\) extends
boundedly in the closed quadratic-form norm, equivalently the graph norm
of the averaged supercharge stack for the stated supersymmetric
realization. This proves \(P_{\sigma}\) Q(H) subset Q(H); preservation
of the Hamiltonian operator domain D(H), or boundedness in its operator
graph norm, is not asserted. Exhaustion by compact slow supports gives
the same identification on an unbounded selected branch, using these
UNIFORM estimates rather than local ellipticity constants that might
depend on \(r_{\max}\).

The restriction of the closed quadratic form of \(H\) to Q(H) intersect
\(\ker (U_{\sigma}^{*})\) is closed and strictly L2-coercive by (1.5)
and (7.3). For \(f\) in the slow form domain, the mixed functional is
continuous in that complementary form norm by (6.1) and the exact source
dual estimate. Riesz representation gives a unique minimizing \(v_{f}\)
in Q(H) intersect \(\ker (U_{\sigma}^{*})\) for each \(f\) in
\(Q(H_{s})\). This is form-domain attainment only; no operator-domain
regularity of \(v_{f}\) is inferred. The resulting Schur form is
continuous in the shifted \(H_{s}\) form norm, and (1.1) makes the
shifted quadratic-form norms equivalent after a bounded scalar shift. It
is therefore closed.

These are form statements. They do not claim a bounded operator \(B\) or
a second-slow-derivative operator identity on arbitrary \(f\), and do
not assume an inverse for \(H_{I}\). All cutoffs, projections, graph
dressings, and estimates intertwine under the same-branch
residual/color/normal-frame unitaries. Different cluster branches and
off-cone selector transport remain outside this local theorem.

\subsection{13. Explicit module and differentiated-tail
checks}\label{i-explicit-module-and-differentiated-tail-checks}

The following coefficient identities make the retained-module and
differentiated-tail steps explicit.

\subsubsection{13.1 Internal potential/Yukawa commute with the actual
injection}\label{i-internal-potentialyukawa-commute-with-the-actual-injection}

At a fixed intrinsic branch, the orthogonal Lie-algebra splitting into
off-block and retained block/center directions gives the graded tensor
product of the fast and slow Clifford modules. The injection is

% DISPLAY i.68
% SOURCE U(A,Z)f = g_(A,Z)(Y) v_F(A,Z) tensor f,
\begin{gather*}
\begin{aligned}
& U(A,Z)f = g_{A,Z}(Y) v_{F}(A,Z) \otimes  f,
\end{aligned}\notag
\end{gather*}

with the identity map on the ENTIRE slow Clifford module. There is no
projection onto an internal singlet or internal eigenstate. Each complex
bifundamental fast mass has \(8 n_{a} n_{b}\) occupied negative modes,
an even number. This remains constant under the gapped continuation, so
the product determinant \(v_{F}\) is even. In particular its graded
tensor injection intertwines the internal odd Clifford generators with
the same retained representation. The internal
Hamiltonian\textquotesingle s Yukawa term is even in any case: it is a
sum of bilinears in retained internal fermions only.

Pointwise in \((A,Z), V_{I}\) is scalar multiplication and \(Y_{I}(A)\)
acts only on those retained Clifford factors. Therefore, as actual maps,

% DISPLAY i.69
% SOURCE V_I U=U V_I,    Y_I U=U Y_I,
% SOURCE U^*Y_I v=Y_I U^*v.
\begin{gather*}
\begin{aligned}
& V_{I} U=U V_{I}, \quad  Y_{I} U=U Y_{I}, \\
& U^{*}Y_{I} v=Y_{I} U^{*}v.
\end{aligned}\notag
\end{gather*}

This proves their zero mixed compression for \(U^{*}v=0\). The cutoff is
a scalar bosonic multiplier, so the identities hold exactly for
\(U_{\sigma}\) as well.

The pair negative-projector radial transport is implemented on the fast
Fock module by a fast-fermion even unitary; its bosonic covariance
transport also acts only on the fast module. Both commute with the
retained Clifford action. The line may carry a nontrivial scalar Berry
connection, but this does not rotate the slow Clifford representation.
Differentiation gives

% DISPLAY i.70
% SOURCE p_A(Uf)=U nabla_A^B f+(1-P)(p_A U)f.
\begin{gather*}
\begin{aligned}
& p_{A}(Uf)=U \nabla _{A}^{B} f+(1-P)(p_{A} U)f.
\end{aligned}\notag
\end{gather*}

The transverse term is retained in (5.7); the Berry term is retained in
\(H_{s}^{B}\) and compared by the full charge estimate (6.5). Nothing in
the radial identification sets the full derivative of \(U\) to zero.
Residual color transformations act simultaneously on the fast
determinant line and the retained module, and \(U\) is the corresponding
equivariant intertwiner. These pointwise tensor identities therefore
descend to the residual-equivariant/physical subspace without requiring
individual block singlets.

\subsubsection{13.2 Exact beta-tail decomposition before internal
differentiation}\label{i-exact-beta-tail-decomposition-before-internal-differentiation}

Use the notation of Section \(2\), and set

% DISPLAY i.71
% SOURCE B=Psi^*F^(-1),
% SOURCE B_1=[C_Y F^(-1)+K^* ; I_Y F^(-1)],
% SOURCE R_m=T^*WT-I.
\begin{gather*}
\begin{aligned}
& B=\Psi ^{*}F^{-1}, \\
& B_{1}=[C_{Y} F^{-1}+K^{*} ; I_{Y} F^{-1}], \\
& R_{m}=T^{*}WT-I.
\end{aligned}\notag
\end{gather*}

The center row of \(B\) is

% DISPLAY i.72
% SOURCE B_C=g_C^(-1/2)[C_Y+K^*(F+L)]F^(-1),
\begin{gather*}
\begin{aligned}
& B_{C}=g_{C}^{-1/2}[C_{Y}+K^{*}(F+L)]F^{-1},
\end{aligned}\notag
\end{gather*}

and the internal row is \(B_{I}=I_{Y} F^{-1}\). Direct substitution into
\(\beta =A^{*}a\) gives the exact identities

% DISPLAY i.73
% SOURCE beta=B T^*WT(d_0+n),      beta_1=B_1d_0,
% SOURCE 
% SOURCE beta-beta_1=(B-B_1)d_0+B R_m d_0+B T^*WT n.           (13.1)
\begin{gather*}
\begin{aligned}
& \beta =B T^{*}WT(d_{0}+n), \quad  \beta _{1}=B_{1} d_{0},
\end{aligned}\notag \\
\begin{aligned}
& \beta -\beta _{1}=(B-B_{1})d_{0}+B R_{m} d_{0}+B T^{*}WT n.
\end{aligned}\tag{13.1}
\end{gather*}

The internal part of \(B-B_{1}\) is zero, and its center part has the
exact factorization

% DISPLAY i.74
% SOURCE (B-B_1)_C d_0
% SOURCE   =(g_C^(-1/2)-I)(C_Y F^(-1)d_0+K^*d_0)
% SOURCE      +g_C^(-1/2)K^* L F^(-1)d_0.                    (13.2)
\begin{gather*}
\begin{aligned}
& (B-B_{1})_{C} d_{0} \\
& =(g_{C}^{-1/2}-I)(C_{Y} F^{-1}d_{0}+K^{*}d_{0}) \\
& +g_{C}^{-1/2}K^{*} L F^{-1}d_{0}.
\end{aligned}\tag{13.2}
\end{gather*}

The metric difference is \(O(\sigma ^{2})\), and the last term has the
explicit \(O(\sigma )\) left factor \(K^{*}\). The other factors in
(13.2) are same-root or input-root-denominator triangle columns.

To expand the middle term of (13.1) without exposing a bare frozen
vacuum norm, use \(T-I=-Z_{m} T\) and the fact that \(T\) commutes with
\(Z_{m}\). Exactly,

% DISPLAY i.75
% SOURCE R_m d_0
% SOURCE   =-T^* Z_m^*d_0 -T^*T Z_m d_0
% SOURCE      +T^*K(K^*d_0)-T^*KK^*T(Z_m d_0).              (13.3)
\begin{gather*}
\begin{aligned}
& R_{m} d_{0} \\
& =-T^{*} Z_{m}^{*}d_{0} -T^{*}T Z_{m} d_{0} \\
& +T^{*}K(K^{*}d_{0})-T^{*}KK^{*}T(Z_{m} d_{0}).
\end{aligned}\tag{13.3}
\end{gather*}

Thus it contains only uniformly bounded left factors multiplying
\(Z_{m}^{*}d_{0}, Z_{m} d_{0}\), or \(K^{*}d_{0}\). These are the
controlled rootwise columns. For example

% DISPLAY i.76
% SOURCE Z_m d_0=F^(-1)L^*d_0
% SOURCE              -F^(-1)C_Y^*[C_Y R^(-1)d_0],
% SOURCE Z_m^*d_0=L F^(-1)d_0-K[C_Y F^(-1)d_0],
% SOURCE K^*d_0=C_Y R^(-1)d_0.
\begin{gather*}
\begin{aligned}
& Z_{m} d_{0}=F^{-1}L^{*}d_{0} \\
& -F^{-1}C_{Y}^{*}[C_{Y} R^{-1}d_{0}], \\
& Z_{m}^{*}d_{0}=L F^{-1}d_{0}-K[C_{Y} F^{-1}d_{0}], \\
& K^{*}d_{0}=C_{Y} R^{-1}d_{0}.
\end{aligned}\notag
\end{gather*}

The last term of (13.1) has the same explicit \(O(\sigma )\) outer \(B\)
multiplying the triangle column \(n\). This proves the extra small
factor by an exact finite identity, not by a formal infinite series.

For an INTERNAL derivative
\(\partial _{j}, Y, K, L, N_{Y}, I_{Y}, C_{Y}, W\) and \(g_{C}\) are
independent of A. Consequently

% DISPLAY i.77
% SOURCE ||partial_j B||<=C sigma/r_*,
% SOURCE partial_j Z_m=-F^(-1)(partial_j F)Z_m,
% SOURCE ||partial_j T||<=C sigma/r_*.
\begin{gather*}
\begin{aligned}
& \Vert \partial _{j} B\Vert \le C \sigma /r_{*}, \\
& \partial _{j} Z_{m}=-F^{-1}(\partial _{j} F)Z_{m}, \\
& \Vert \partial _{j} T\Vert \le C \sigma /r_{*}.
\end{aligned}\notag
\end{gather*}

Every derivative of a root-preserving coefficient \(F_{e}^{-1}\) adds
\(1/r_{e}\). A derivative of \(F_{e}^{-1}C_{e}^{*}\) can remove its
optional \(O(\kappa )\) factor, but replaces it by \(O(1/r_{e})\). The
same-root or triangle column therefore still has an extra
\(1/r_{e}\le 1/r_{*}\) in its low Gaussian norm. Crucially the outer
\(B, K^{*}\), or \(g_{C}^{-1/2}-I\) displayed in (13.1)-(13.3) remains,
or its derivative has the SAME \(\sigma\) gain and another \(1/r_{*}\).

Applying the finite Gaussian rootwise bounds to these differentiated
identities gives

% DISPLAY i.78
% SOURCE ||(beta-beta_1)U||<=C sigma W_2^(1/2),
% SOURCE ||[partial_A(beta-beta_1)]U||
% SOURCE                     <=C sigma r_*^(-1)W_2^(1/2).    (13.4)
\begin{gather*}
\begin{aligned}
& \Vert (\beta -\beta _{1})U\Vert \le C \sigma  W_{2}^{1/2}, \\
& \Vert [\partial _{A}(\beta -\beta _{1})]U\Vert \\
& \le C \sigma  r_{*}^{-1}W_{2}^{1/2}.
\end{aligned}\tag{13.4}
\end{gather*}

The derivative in the second line acts on the operator coefficient,
exactly as required by the curvature term, not on \(U\). Derivatives of
\(U\) are already present in the complete-charge cross. At
\(A=0, \partial _{A} d_{0}\) need not vanish; (13.1)-(13.3) explicitly
retain its contribution with the same extra fast factor. The non-small
derivative of \(\beta _{1}\) itself is not included in the tail estimate
and remains in the leading H2 scalar.

Finally the coefficients of \(\beta -\beta _{1}\) have at least two
fast-coordinate factors before \(p_{F}\). A fast divergence removes at
most one. Differentiating the displayed normalized coefficients
therefore gives
\(\vert \operatorname{div} _{F}(b-b_{1})\vert \le C \sigma /r_{*}\); an
additional internal derivative costs at most \(1/r_{*}\) and leaves this
\(\sigma\) factor. This proves the divergence portions of Section
\(4.3\) with the required gain as well.

\subsubsection{13.3 Form domain
convention}\label{i-form-domain-convention}

The infimum, projection bound, and Riesz minimizer are statements about
Q(H), the closed quadratic-form domain. They do not establish invariance
of the Hamiltonian operator domain D(H). The minimizing complement is
asserted only in Q(H), for \(f\) in \(Q(H_{s})\). The global lower bound
does not require attainment.
